% Generated draft: requires mathematical and source review.
\documentclass[11pt]{article}
\usepackage[margin=25mm]{geometry}
\usepackage{fontspec}
\setmainfont{Latin Modern Roman}
\usepackage{amsmath,amssymb,mathtools,unicode-math}
\setmathfont{Latin Modern Math}
\usepackage{enumitem,needspace,xurl,hyperref}
\setlength{\parindent}{0pt}
\setlength{\parskip}{5pt}
\setlength{\emergencystretch}{3em}
\newcommand{\N}{\mathbb{N}}
\newcommand{\Z}{\mathbb{Z}}
\newcommand{\Q}{\mathbb{Q}}
\newcommand{\R}{\mathbb{R}}
\newcommand{\C}{\mathbb{C}}
\newcommand{\F}{\mathbb{F}}
\newcommand{\A}{\mathbb{A}}
\newcommand{\PP}{\mathbb{P}}
\newcommand{\E}{\mathbb{E}}
\newcommand{\eps}{\varepsilon}
\newcommand{\dd}{\,\mathrm{d}}
\begin{document}
\section*{UnsolvedMath ID: 30001890}
\textbf{Source problem code:} catalogue entry
\subsection{Definitions and mathematical statement}
Let \(H\) be a family of forbidden geometric configurations, and let \(n\) be an admissible integer vertex count. The source does not specify the geometric graph model or the meaning of containment of a configuration. Accordingly, the following notation is conditional on fixing those conventions.

Write \(\mathcal G_n\) for the class of geometric graphs with exactly \(n\) vertices in the chosen model, and write \(F\preccurlyeq G\) when the configuration \(F\) occurs in \(G\) according to the chosen containment relation. A graph \(G\) is \(H\)-free if
\[
\forall F\in H,\qquad F\not\preccurlyeq G.
\]
Define the set of attainable edge counts by
\[
A_H(n)=\bigl\{|E(G)|:G\in\mathcal G_n,\ \forall F\in H,\ F\not\preccurlyeq G\bigr\}.
\]
\textbf{Problem.} Given \(H\), determine the maximum number of edges in an \(H\)-free geometric graph with \(n\) vertices; in the notation above, determine
\[
\operatorname{ex}_{\mathcal G,\preccurlyeq}(n,H)=\max A_H(n)
\]
whenever this maximum exists.

This is a conditional formulation of the supplied question. No particular family \(H\), ambient space, drawing convention, general-position hypothesis, or containment relation is supplied. The source also gives no convention for an empty admissible class or for a failure of a finite maximum to exist. These choices must be resolved before the question defines a particular extremal function.
\subsection{Short English statement}
For a given forbidden family \(H\), determine how many edges an \(n\)-vertex geometric graph can have while avoiding every configuration in \(H\). The geometric model and the meaning of occurrence must first be specified; the supplied record leaves them unresolved.
\subsection{Sources}
\begin{itemize}
\item Supplied local UnsolvedMath registry record: permanent ID \(30001890\), title ``Extremal Geometric Graphs with Forbidden Configurations,'' source problem code \texttt{catalogue entry}, in the Oberwolfach Reports Open Problems collection. No individual proposer is supplied; the proposed year is recorded as 2011.
\item Supplied source citation: \emph{Discrete Geometry (2012)}, Oberwolfach Reports, pp.\ 2459--2548, DOI \texttt{10.4171/owr/2011/44}. Supplied report link and literature-source link, labeled ``OWR 44 (2012)'': \url{https://doi.org/10.4171/owr/2011/44}. The record describes Oberwolfach Reports as documenting mathematical workshops and their open-problem discussions; it supplies no exact problem page or individual attribution.
\item Literature assessment dated exactly \texttt{2026-08-21}: the supplied curation labels the problem open and states that the source frames the extremal geometric-graph problem for general forbidden configurations, with no general asymptotic formula matching the extracted question located. This is supplied dated literature triage, not an independently verified assessment; time-sensitive status and cited claims require independent verification.
\item The registry labels the statement \texttt{corrected\_verified}, but its supplied verification explains only that math delimiters were added around the forbidden family and vertex-count parameter without changing the question. This does not establish mathematical completeness or independent source verification.
\end{itemize}
\end{document}
